\documentclass[10pt]{amsart}

\usepackage[a4paper,margin=0.82in]{geometry}
\usepackage{amsmath,amssymb,mathtools}
\usepackage[hidelinks]{hyperref}

\newtheorem{theorem}{Theorem}
\newtheorem{corollary}[theorem]{Corollary}
\newtheorem{lemma}[theorem]{Lemma}
\theoremstyle{remark}
\newtheorem{remark}[theorem]{Remark}

\newcommand{\ex}{\operatorname{ex}}
\newcommand{\ind}{\mathbf 1}
\newcommand{\ldir}{\ell^{\mathrm{dir}}}
\newcommand{\kdir}{k^{\mathrm{dir}}}

\title[A sharper directed-connectivity bound]
{A Two-Path Improvement for Directed Connectivity Bounds}
\author{}
\date{}

\begin{document}

\begin{abstract}
For a simple digraph on $n$ vertices in which every local arc-connectivity,
or every local vertex-connectivity, is at most $m-1$, the bound
$\lfloor n^2/4\rfloor+4(m-1)(n-1)$ was previously known.
We observe that such a digraph avoids $m$ directed paths of length two with
common endpoints.  A theorem of Huang and Lyu therefore improves the upper
bound, for every fixed $m\ge 3$ and all sufficiently large $n$, to
\[
  \frac{n^2}{4}+(m-1)n+O_m(1).
\]
Thus the coefficient of the linear error term drops by a factor of four.
\end{abstract}

\maketitle

\section{The two-path theorem}

All digraphs below are \emph{strict}: they have no loops or parallel arcs,
but a pair of opposite arcs is allowed.  For distinct vertices $u,v$ of a
digraph $D$, write
\[
  p_2^D(u,v)
  :=\bigl|\{x\in V(D): ux,xv\in A(D)\}\bigr|.
\]
Thus $p_2^D(u,v)$ is the number of directed $u$--$v$ paths of length two.
Let $P_{s,2}$ denote the union of $s$ such paths with common initial and
terminal vertices and distinct internal vertices.

For an integer $t\ge2$, put
\begin{align*}
 N(t)&:=\max\left\{
 t^3+4t^2+3t+4,\,
 \frac{17}{2}t^2+30t+27
 \right\},\\
 g(n,t)&:=
 \left\lceil\frac{n+t}{2}\right\rceil
 \left\lfloor\frac{n-t}{2}\right\rfloor
 +tn+1.
\end{align*}

We use the following result of Huang and Lyu.

\begin{theorem}[Huang--Lyu \cite{HuangLyu}]
\label{thm:huang-lyu}
Let $t\ge2$ and $n\ge N(t)$.  Then
\[
  \ex(n,P_{t+1,2})\le g(n,t).
\]
More precisely, equality holds when
$\lfloor(n-t)/2\rfloor$ is odd, while for even
$\lfloor(n-t)/2\rfloor$ the value belongs to
$\{g(n,t)-1,g(n,t)\}$.
\end{theorem}

The form needed here is an immediate consequence.

\begin{lemma}[Improved counting core]
\label{lem:core}
Let $C\ge2$ and $n\ge N(C)$.  Suppose that a strict digraph $D$ satisfies
\begin{equation}
\label{eq:budget}
  \ind_{uv\in A(D)}+p_2^D(u,v)\le C
  \qquad\text{for all distinct }u,v\in V(D).
\end{equation}
Then
\begin{equation}
\label{eq:g-bound}
 |A(D)|\le g(n,C).
\end{equation}
In particular,
\begin{equation}
\label{eq:asym-core}
 |A(D)|
 \le \frac{n^2}{4}+Cn+1-\frac{C^2}{4}
 =\frac{n^2}{4}+Cn+O_C(1).
\end{equation}
\end{lemma}

\begin{proof}
Condition \eqref{eq:budget} implies $p_2^D(u,v)\le C$ for every ordered
pair $(u,v)$.  Hence $D$ contains no copy of $P_{C+1,2}$, and
Theorem~\ref{thm:huang-lyu} gives \eqref{eq:g-bound}.

For the displayed estimate, observe that
\[
 \left\lceil\frac{n+C}{2}\right\rceil
 \left\lfloor\frac{n-C}{2}\right\rfloor
 \le \frac{n^2-C^2}{4}.
\]
Substitution in the definition of $g(n,C)$ proves \eqref{eq:asym-core}.
\end{proof}

\section{Application to bounded local connectivity}

For a strict digraph $D$, let $\lambda_D(u,v)$ be the maximum number of
arc-disjoint directed $u$--$v$ paths, and let $\kappa_D(u,v)$ be the maximum
number of internally vertex-disjoint directed $u$--$v$ paths.  Define
$\ldir_m(n)$ and $\kdir_m(n)$ to be the largest number of arcs in an
$n$-vertex strict digraph satisfying, respectively,
\[
 \lambda_D(u,v)\le m-1
 \quad\text{or}\quad
 \kappa_D(u,v)\le m-1
 \qquad (u\ne v).
\]

\begin{theorem}[Sharper linear error term]
\label{thm:main}
Let $m\ge3$ and $n\ge N(m-1)$.  Then
\begin{equation}
\label{eq:main-exact}
 \ldir_m(n)\le \kdir_m(n)
 \le
 \left\lceil\frac{n+m-1}{2}\right\rceil
 \left\lfloor\frac{n-m+1}{2}\right\rfloor
 +(m-1)n+1.
\end{equation}
Consequently,
\begin{equation}
\label{eq:main-asym}
 \ldir_m(n),\ \kdir_m(n)
 \le
 \frac{n^2}{4}+(m-1)n+O_m(1).
\end{equation}
\end{theorem}

\begin{proof}
Fix distinct vertices $u,v$.  The direct arc $uv$, when present, together
with all directed paths $u\to x\to v$, is a pairwise arc-disjoint family:
different middle vertices use different first and last arcs.  It is also a
pairwise internally vertex-disjoint family, since the two-step paths have
distinct internal vertices and the direct path has none.  Therefore either
of the hypotheses
\[
 \lambda_D(u,v)\le m-1,
 \qquad
 \kappa_D(u,v)\le m-1
\]
implies
\[
 \ind_{uv\in A(D)}+p_2^D(u,v)\le m-1.
\]
Lemma~\ref{lem:core}, with $C=m-1$, now gives the upper bound in
\eqref{eq:main-exact}.  Finally, every arc-feasible digraph is
vertex-feasible because $\kappa_D(u,v)\le\lambda_D(u,v)$, so
$\ldir_m(n)\le\kdir_m(n)$.  Equation \eqref{eq:main-asym} follows from
\eqref{eq:asym-core}.
\end{proof}

The exact parity form of the new upper bound is occasionally convenient.
With $t=m-1$,
\[
 g(n,t)=\frac{n^2}{4}+tn+1-
 \begin{cases}
   t^2/4,&n\equiv t\pmod 2,\\[1mm]
   (t+1)^2/4,&n\not\equiv t\pmod 2.
 \end{cases}
\]

\begin{corollary}[Comparison with the conjectured value]
\label{cor:sandwich}
For every fixed $m\ge3$ and all sufficiently large $n$,
\begin{equation}
\label{eq:sandwich}
 \left\lfloor\frac{(n+m-2)^2}{4}\right\rfloor
 \le \ldir_m(n)\le\kdir_m(n)
 \le g(n,m-1).
\end{equation}
Hence
\[
 \frac{n^2}{4}+\frac{m-2}{2}n+O_m(1)
 \le \ldir_m(n)\le\kdir_m(n)
 \le \frac{n^2}{4}+(m-1)n+O_m(1).
\]
\end{corollary}

\begin{proof}
The lower bound is the standard augmented bipartite construction: take all
arcs from a part $A$ to a part $B$, and give every vertex of $B$ exactly
$m-2$ cyclic in-neighbours inside $B$, with
$|B|=\lceil(n+m-2)/2\rceil$.  It has
$\lfloor(n+m-2)^2/4\rfloor$ arcs and local arc-connectivity at most $m-1$.
The remaining inequalities are Theorem~\ref{thm:main}.
\end{proof}

\begin{remark}
Theorem~\ref{thm:main} improves the coefficient $4(m-1)$ in the previously
known bound
\[
 \left\lfloor\frac{n^2}{4}\right\rfloor+4(m-1)(n-1)
\]
to $m-1$, at the cost of requiring $n\ge N(m-1)$; compare
\cite[Theorems 1.11 and 1.12]{Vandenbruwaene}.  It does not prove the exact
conjecture.  The condition of being $P_{m,2}$-free remembers only paths of
length two, whereas bounded local connectivity also restricts the direct arc
and all longer paths.  That discarded information is precisely where a
further improvement would have to enter.
\end{remark}

\begin{thebibliography}{9}

\bibitem{HuangLyu}
Z.~Huang and Z.~Lyu,
\emph{Extremal digraphs containing at most $t$ paths of length $2$ with the
same endpoints},
Discrete Appl. Math. \textbf{388} (2026), 1--10.

\bibitem{Vandenbruwaene}
L.~Vandenbruwaene,
\emph{Maximising Edge Density Subject to Connectivity Constraints:
Erd\H{o}s Problem 915, from Proof to Search},
Master's thesis, KU Leuven, 2026.

\end{thebibliography}

\end{document}
